\newpage
\setcounter{page}{1}
\thispagestyle{empty}   
\begin{center}
\begin{minipage}{\textwidth}
    \centering
    \large
    \MakeUppercase{\textbf{
     \vspace{1cm}
     \noindent\rule{\textwidth}{0.4pt} 
         \titlename \\} }
         \authorShort \\
        \departmentname \\
        \studentID@nmct.edu.mn
\end{minipage}
\thispagestyle{empty}

\noindent\textbf{Хураангуй:}
{\hspace{135mm}}    
\vspace{3mm}
\begin{flushleft}
\hspace{1em}13-р зууны Монголын эзэнт гүрний түүхийг (Чингис хаан, байлдан дагуулалт, соёл)
интерактив мобайл апп-аар таниулах төсөл нь залуу үеийнхэнд түүхийн мэдлэгийг түгээх, сонирхлыг нэмэгдүүлэхэд
чиглэсэн. Апп нь Flutter-ээр хөгжүүлэгдэж, оффлайн горимд (SQLite) ажиллах бөгөөд онлайн болоход Firebase-ээр sync 
хийнэ. Гол функцууд: түүхэн хүмүүсийн намтар, интерактив газрын зураг, цагийн хэлхээс, квиз, мультимедиа. Дизайн нь модульчлагдсан 
(UI, Data, Logic), аюулгүй (AES шифрлэлт, биометр баталгаажуулалт). Төсөл Agile аргачлалаар 8 сарын хугацаатай хэрэгжинэ. 
UML диаграммууд (Use Case, Sequence, System) болон SRS, SDD баримт бичгүүд бэлэн. Үр дүн: Монгол хэл дээрх боловсролын
дижитал нөөцийг баяжуулах, түүхийн мэдлэгийг хүртээмжтэй болгох.
\end{flushleft}

\vspace{3mm}

\noindent\textbf{Түлхүүр үгс:} 13-р зуун, Монголын түүх, Интерактив апп,Flutter хөгжүүлэл
\end{center}